\plainsection{摘要}

\vspace{0.45\baselineskip}

当两支队伍共用一条只能单车通过的狭窄通道时，前车是否清空、当前任务何时可以开始、空中观察是否仍然有效，都会影响下一辆车的发车时间。本设计在 Webots 中搭建了空地协同对抗靶场，由六辆差速地面车和两架无人机组成红蓝两队，登记运输、侦察、报告和巡察等 14 项任务。场地设置入口泊位、交付点和外侧路线，裁判根据车辆到位、驻留、通道占用和任务期限记录各回合的结果。

地面车利用定位和航向反馈控制左右轮速，在通道外完成对准，到达交付点后减速驻留。放行规则按车体包络判断前车是否清空。蓝方无人机提供带采样时刻的观察报告，第二批运输收到有效观察后发车，首帧过期时可请求补采。协同设计参考群集避让、响应阈值分工和知情个体引导，并在离散算例中借鉴军蚁搭桥分析支援收益与新增车时的关系。

正常、首帧延迟与侦察车起点变化三个基础回合均为 14/14 项按期。首帧延迟到过期时，主动补采使按期数从 13 项恢复到 14 项。11 月 20 日的 AMAX 补验沿用同一任务卡。低箱整场中，B1 根据低位雷达回波更新局部地图，在 133.776~s 绕行交付，整场零碰箱。双机旋翼整场中，B2 释放任务，B1 接手并使用 B3 新拍的相机帧放行，在 340.176~s 交付。两轮均为 14/14 项按期，通道 5 进 5 清、零占用重叠。观察链路持续断开时，B2 保持等待，其余 13 项按期。实验结论对应指定布局、固定中转点与软件故障注入。

\noindent{\bfseries 关键词}\quad 智能靶场、空地协同、共享通道、任务时间窗、Webots
